\section{Introduction}
\label{sec:intro}

Blockchain sharding is the dominant approach to scaling permissionless consensus without a
linear increase in per-node work~\cite{luu2016elastico,kokoris2018omniledger,zamani2018rapidchain}.
Rather than asking every validator to validate every transaction, sharded designs partition the
validator set into \emph{committees}, each of which runs an independent BFT consensus instance
over a slice of state, and stitch the slices together with a cross-shard transaction protocol.
The security argument for an individual committee is the standard BFT argument: with $N=3F+1$
validators and a commit quorum of $2F+1$, any two commit quorums intersect in at least $F+1$
validators, so two conflicting commits require at least $F+1$ validators to \emph{equivocate}
(sign both branches). As long as fewer than $F+1$ validators are adversarial, shard safety holds.

\begin{figure}[t]
\centering
\begin{tikzpicture}[
  font=\scriptsize,
  >={Latex[length=1.3mm]},
  box/.style={rounded corners=2pt, draw=black!55, fill=black!4, align=center,
              inner sep=2.5pt, text width=0.84\columnwidth, minimum height=6mm},
  flow/.style={->, line width=.5pt, black!65},
]
% ---------- (a) the attack (vertical pipeline) ----------
\node[box] (n1) {Global validator set \textcolor{red!70!black}{(adversary $<\tfrac13$)}};
\node[box, below=3mm of n1] (n2) {Committee $N{=}3F{+}1$: \textcolor{red!70!black}{bribe $F{+}1$ equivocators}};
\node[box, below=3mm of n2] (n3) {Two conflicting certificates $\sigma_A\!\neq\!\sigma_B$};
\node[box, below=3mm of n3] (n4) {Extract value $V$ (cross-shard / bridge)};
\draw[flow] (n1)--(n2); \draw[flow] (n2)--(n3); \draw[flow] (n3)--(n4);
\node[anchor=west, font=\scriptsize\bfseries] at ([yshift=3mm]n1.north west) {(a)~The attack};

% ---------- (b) conventional vs.\ ours: the accountability race ----------
\begin{scope}[shift={($(n4.south west)+(0,-0.65)$)}]
  \node[anchor=west, font=\scriptsize\bfseries] at (0,0)
     {(b)~Race: settlement $\Tsettle$ vs.\ accountability $\Tacc$};
  % conventional (ungated): settlement first -> attack pays
  \node[anchor=west] at (0,-0.82) {Conventional:};
  \draw[->, black!55] (1.65,-0.82)--(6.15,-0.82);
  \filldraw[red!70!black] (2.55,-0.82) circle (1.2pt) node[above=2pt, red!70!black] {$\Tsettle$};
  \filldraw[black!55]     (4.75,-0.82) circle (1.2pt) node[above=2pt, black!60] {$\Tacc$};
  \node[anchor=west] at (6.3,-0.82) {\textcolor{red!70!black}{\ding{55}}};
  \node[anchor=west] at (1.65,-1.32) {\textcolor{red!70!black}{$\Tsettle{<}\Tacc$}: adversary keeps $V$};
  % this paper (finality-gated): accountability first -> attack fails
  \node[anchor=west] at (0,-2.25) {This paper:};
  \draw[->, black!55] (1.65,-2.25)--(6.15,-2.25);
  \filldraw[black!55] (3.05,-2.25) circle (1.2pt) node[above=2pt, black!60] {$\Tacc$};
  \draw[green!60!black, line width=.9pt] (4.0,-2.47)--(4.0,-2.03);
  \filldraw[green!55!black] (5.15,-2.25) circle (1.2pt) node[above=2pt, green!50!black] {$\Tsettle$};
  \node[anchor=west] at (6.3,-2.25) {\textcolor{green!55!black}{\ding{51}}};
  \node[anchor=west] at (1.65,-2.75) {\textcolor{green!55!black}{$\Tsettle{\ge}\Tacc$}: $\pwin{=}0$ (finality gate)};
\end{scope}
\end{tikzpicture}
\caption{Overview. \textbf{(a)}~An adversary below the global $\tfrac13$ threshold bribes the $F{+}1$
validators quorum intersection makes \emph{necessary} in one $N{=}3F{+}1$ committee, forging two
conflicting certificates and moving value $V$ across the shard/bridge boundary. \textbf{(b)}~The attack
pays only if value settles before accountability: conventional ungated settlement gives
$\Tsettle{<}\Tacc$ (adversary keeps $V$), whereas gating cross-shard acceptance on finality forces
$\Tsettle{\ge}\Tacc$ and $\pwin{=}0$ (Proposition~\ref{prop:immunity}).}
\label{fig:overview}
\end{figure}

This per-committee argument hides a well-known tension. Because voting power is split, the
\emph{absolute} number of validators an adversary must control to violate safety in one shard is
far smaller than the number required to attack a monolithic chain---a single corrupt committee
``can compromise all the transactions that touch'' it, and detection and recovery remain open
problems~\cite{bano2019sok}. In parallel, a maturing literature shows that validator faults need
not be obtained by intrusion at all: they can be \emph{bought}. Smart contracts can bribe miners
and validators trustlessly~\cite{mccorry2018smart,winzer2019temporary,tran2023crosschain},
game-theoretic analyses price bribery against proof-of-stake \emph{safety}~\cite{karakostas2024bribing},
and recent work implements and evaluates trustless contracts that buy votes, induce validator
exits, and auction randomness-beacon manipulation~\cite{sookitoth2025bribers,kelkar2025omerta,austgen2025liquefaction}.

\paragraph{The composition is not, by itself, an attack.}
Taken together these two strands appear to compose into a devastating attack: pay $F+1$
validators in a small committee to equivocate, and break shard safety for the price of a handful
of bribes. Yet the composition is \emph{not} obviously an attack with consequence. Accountable
BFT consensus turns equivocation into a transferable cryptographic proof and slashes the offending
stake~\cite{buterin2017casper,buchman2018tendermint}. If the adversary forges two conflicting
certificates but the system slashes the bribed validators and reverts the bad branch before any
value moves, the adversary has merely burned the validators' stake and its own bribes for nothing.
The attack's profitability therefore hinges on a question the prior literature does not answer in
the sharded setting: \emph{can the adversary realize value before accountability catches up?}

\paragraph{The accountability race.}
This paper studies that question directly. Our central object is the \emph{accountability race}
(Figure~\ref{fig:overview}) between the time to convert a fraudulently finalized state into irreversible
value---$\Tsettle$, dominated by cross-shard receipt processing and bridge withdrawal beyond
clawback---and the time for the protocol to observe the conflicting certificates, propagate the
equivocation proof, slash, and revert---$\Tacc$. We argue that this race, not the raw cost of
corruption, is what determines whether sub-threshold committee bribery is a real threat to a
given sharded design, and we give the conditions under which it is. Sharding changes the
\emph{economics} of BFT safety: an adversary can buy a \emph{local} shard safety violation after
committee assignment, instead of controlling a large network-wide Byzantine fraction---but the
violation pays only under precise timing and economic conditions, which we make explicit and
checkable.

\paragraph{Two clarifications we treat as load-bearing.}
First, ``sub-threshold'' means below the \emph{network-wide} corruption threshold, \emph{not}
below a committee's local BFT safety threshold: the adversary is globally small but locally
decisive (\S\ref{sec:threat}). Second, the $F+1$ equivocators are \emph{necessary} but not
\emph{sufficient}. By quorum intersection two conflicting $2F{+}1$ certificates must overlap in at
least $F{+}1$ equivocators (Lemma~\ref{lem:intersection}); but each certificate still needs $F$
honest signatures, so the attack also requires honest validators to be \emph{split} across the
two branches---through leader equivocation, message delay, transient asynchrony, network
partition, or protocol-specific scheduling. We fold the probability that a usable split
materializes into a success probability $\psucc$ and never claim that buying $F+1$ equivocations
alone forces a safety break.

\paragraph{Contributions.}
\begin{itemize}
  \item \textbf{A threat model for sub-threshold committee bribery} (\S\ref{sec:threat}) that
  makes precise the sense in which an adversary is globally sub-threshold yet locally decisive,
  and separates the \emph{necessary} purchase of $F+1$ equivocations from the \emph{sufficient}
  honest-split condition captured by $\psucc$.
  \item \textbf{A bribe-cost model} (\S\ref{sec:cost}) in which the equilibrium price of
  equivocation is $b_i \gtrsim q\,(s_i + R_i)$, where $q$ is the probability the fraud proof is
  enforced before the validator's stake becomes unslashable. This connects the cost of corruption
  to accountability latency rather than treating the Byzantine fraction as exogenous.
  \item \textbf{An accountability-race model} (\S\ref{sec:race}) factoring attacker profit into a
  \emph{window} gate $\Tsettle<\Tacc$ and an \emph{economic} gate $V > B + \Cop$, with $\Tacc$ a
  master variable that couples them, and an \emph{immunity-by-construction} result
  (Proposition~\ref{prop:immunity}): finality-gating cross-shard acceptance forces $\pwin=0$.
  \item \textbf{A reconfiguration analysis} (\S\ref{sec:reconf}) that extends the model to
  epoch-boundary state and receipt handoff. We prove a reconfiguration-safe-handoff theorem
  (Proposition~\ref{prop:reconf-safe}) and a residual-vulnerability theorem
  (Proposition~\ref{prop:reconf-residual}), give a taxonomy of six handoff designs, and four
  explicit counterexamples that bound when the attack does \emph{not} apply.
  \item \textbf{A reproducible artifact with a measured testbed and a real-systems case study}
  (\S\ref{sec:impl}--\S\ref{sec:eval}): a seeded simulator, a local proof-of-concept that produces
  genuine equivocation evidence and drives a model \textsc{PayToEquivocate} escrow, and a
  discrete-event BFT testbed with real Ed25519 signatures in which $\Tacc$ and $\Tsettle$ are
  \emph{measured} from a running protocol---independently reproducing the model's $\rho\approx1$
  threshold (the harness also reproduces the closed-form success probability to within Monte-Carlo
  error, max.\ deviation $\ArtifactMaxValErr$). A case study on six deployed systems' documented
  parameters shows that fast evidence ($q\to1$) and finality-gating place them in the safe regime,
  concentrating residual risk in fast-exit, ungated-settlement designs. Every figure and table is
  regenerated from one configuration file.
\end{itemize}

\paragraph{Positioning.}
We do not claim to invent bribery, smart-contract bribery, the observation that small committees
are cheap to corrupt, or cross-shard/bridge value extraction; all are prior art surveyed in
\S\ref{sec:related}. Prior work studies validator bribery, smart-contract bribery, adaptive shard
corruption, and cross-shard attacks \emph{mostly separately}. Our contribution is their
\emph{composition}: post-assignment conditional bribery of a committee-local quorum-intersection
set, \emph{priced by accountability latency}, with profit determined by whether cross-shard or
bridge settlement outruns slashing---together with explicit viability and immunity conditions and
the reconfiguration extension.
